\scanpage{121}
% archive sheet number: 120
la condition (83) a) s'exprime par
\[
 \sum_i\varphi_i^\omega(\delta)=0 \tag*{(86)}
\]
la condition b) par la condition
\[
 k.\big((\varphi_i^\omega(\delta))_i\big)=F_t \tag*{(87)}
\]
qui détermine d'emblée \deleted{les} $(\varphi_i^\omega(\delta))_i$ à scalaire inv. près -- la condition c les déterminant enfin (pour \uuline{$\unsure{\mathbb{i}|\boldsymbol{i}|i|\dot{\imath}}$} également choisi) au \uline{signe} près.

Bien sûr, \deleted{\unsure{on aura}} \inserted{la relation entre} $\kappa$ et les $\varphi_i^\omega(\delta)$ s'explicite par
\[
 (\varphi_i^\omega(\delta))=\underbrace{\check\kappa(\delta^{\unsure{\vee*|\vee}})}_{\mathcal O_t(-1)}\wedge\omega
 \ \in\mathcal O_t(-1)(\omega)\simeq F_t\subset V_{\unsure{I|J}}(k)\subset k^I \tag*{(88)}
\]
et la condition c) s'écrit Dieu sait comment -- il faudra \deleted{$\varphi_i^\omega(\delta)$} y revenir. Reste à exprimer la dépendance p.r. à $\omega$, pour \deleted{\unsure{exp|expr}} trouver l'invariance de l'épinglage de Legendre p.r. au choix de $\mathbb{i}$ i.e. p.r. au choix de $\omega=\omega_{\unsure{i|\mathbb{i}|\dot{\imath}}}$ ($\alpha:\mu_4^*(k)\simeq\underline\omega$ étant supposé déterminé), on trouve la relation nécessaire
\[
 \varphi_i^{\omega'}(\delta)=\deleted{\unsure{\mathbb{i}_\omega\alpha}}\mathbb{i}_\omega\varphi_i^\omega(\delta) \tag*{(89)}
\]
où
\[
 \mathbb{i}_\omega\overset{\text{déf}}=\alpha(\omega)\qquad
 \begin{aligned}
 &\alpha:\mu_4^*\simeq\underline\omega\text{ étant l'iso.}\\
 &\text{(supposé déterminé)}\\
 &\text{associé à l'épinglage}\\
 &\text{de Legendre cherché}
 \end{aligned}
\]
Posant
